% Part: first-order-logic
% Chapter: completeness
% Section: introduction

\documentclass[../../../include/open-logic-section]{subfiles}

\begin{document}

\iftag{FOL}
      {\olfileid[pt]{fol}{com}{int}}
      {\olfileid[pt]{pl}{com}{int}}

\olsection{Introdução}

O teorema da completude é um dos resultados mais fundamentais sobre
lógica.  Ele vem em duas formulações, cuja equivalência iremos
provar.  Em sua primeira formulação, ele diz algo fundamental sobre
a relação entre consequência semântica e nosso sistema de
!!{derivation}: se !!a{sentence}~$!A$ segue de algumas !!{sentence}s
$\Gamma$, então também existe !!a{derivation} que estabelece $\Gamma
\Proves !A$. Assim, o sistema de !!{derivation} é tão forte quanto pode
possivelmente ser sem provar coisas que de fato não se seguem.

Em sua segunda formulação, ele pode ser enunciado como um resultado de
existência de modelo: todo conjunto consistente de !!{sentence}s é
satisfatível.  A consistência é uma noção da teoria da prova: ela diz
que nosso sistema de !!{derivation} é incapaz de produzir certas
!!{derivation}s.  Mas quem garante que, só porque não há !!{derivation}s
de um certo tipo a partir de~$\Gamma$, está garantido que existe
\iftag{FOL}{!!a{structure}~$\Struct{M}$}{!!{valuation}{~$\pAssign{v}$}
com $\iftag{FOL}{\Sat{M}}{\pSat{v}}{\Gamma}$}? Antes de
o teorema da completude ser provado pela primeira vez --- na verdade,
antes de termos os sistemas de !!{derivation} que temos hoje --- o grande
matemático alemão David Hilbert defendia a visão de que a consistência
das teorias matemáticas garante a existência dos objetos de que elas
tratam. Ele colocou isso da seguinte forma em uma carta a Gottlob
Frege:
\begin{quote}
  Se os axiomas dados arbitrariamente não se contradizem uns aos
  outros com todas as suas consequências, então eles são verdadeiros e
  as coisas definidas pelos axiomas existem. Este é para mim o critério
  de verdade e existência. 
\end{quote}
Frege discordou veementemente. A segunda formulação do teorema da
completude mostra que Hilbert estava certo pelo menos no sentido de que,
se os axiomas são consistentes, então existe \emph{alguma}
\iftag{FOL}{!!{structure}}{!!{valuation}} que os torna todos
verdadeiros.

Estas não são as únicas razões pelas quais o teorema da completude ---
ou melhor, sua prova --- é importante.  Ele tem uma série de
consequências importantes, algumas das quais discutiremos
separadamente.  Por exemplo, como toda !!{derivation} que mostra
$\Gamma \Proves !A$ é finita e, portanto, só pode usar uma quantidade
finita das !!{sentence}s em~$\Gamma$, segue do teorema da completude
que, se $!A$ é uma consequência de~$\Gamma$, então já é uma
consequência de um subconjunto finito de~$\Gamma$.  Isso é chamado de
\emph{compacidade}.  Equivalentemente, se todo subconjunto finito de
$\Gamma$ é consistente, então o próprio $\Gamma$ deve ser consistente.

Embora o teorema da compacidade decorra do teorema da completude
pelo desvio através de !!{derivation}s, também é possível usar a
\emph{prova do} teorema da completude para estabelecê-lo
diretamente. Pois o que a prova faz é tomar um conjunto de !!{sentence}s
com uma certa propriedade --- a consistência --- e construir
!!a{structure} a partir desse conjunto que possui certas propriedades
(neste caso, que ela satisfaz o conjunto). Quase exatamente a mesma
construção pode ser usada para estabelecer a compacidade diretamente,
partindo de conjuntos ``finitamente satisfatíveis'' de !!{sentence}s em
vez de conjuntos consistentes.
\iftag{FOL}{A construção também produz outras consequências, por exemplo,
  que todo conjunto satisfatível de !!{sentence}s tem um modelo finito ou
  !!{denumerable}s.  (Esse resultado é chamado de
  teorema de L\"owenheim--Skolem.)  Em geral, a construção de
  !!{structure}s a partir de conjuntos de !!{sentence}s é usada com frequência
  em lógica, e às vezes até em filosofia.}{}

\end{document}
